\documentclass[11pt]{article}
\usepackage[a4paper,margin=1in]{geometry}
\usepackage{graphicx}
\usepackage{tabularx}
\usepackage{amsmath}
\usepackage{amsfonts}
\usepackage{amssymb}
\usepackage{enumitem}
\usepackage{gensymb}
\usepackage{xcolor}
\usepackage[T1]{fontenc}
\usepackage[utf8]{inputenc}

\usepackage[hidelinks]{hyperref}
\usepackage[
	sorting=none,
	backend=biber,
	hyperref=true,
	natbib=true,
    url=false, 
    doi=true
]{biblatex}
\usepackage{url}
\addbibresource{sources.bib}

\usepackage{lmodern}
\usepackage{mathpazo}
\DeclareUnicodeCharacter{2082}{\textsubscript{2}}
\newcommand{\cotwo}{\text{CO\textsubscript{2}}}

\renewcommand{\thefootnote}{\alph{footnote}}
\renewcommand{\thebibliography}{}
\renewcommand{\textcite}[1]{\citeauthor{#1} (\citeyear{#1})}
\renewcommand{\parencite}[1]{(\citeauthor{#1}, \citeyear{#1})}

\newcommand{\draftmsg}[1]{\textbf{\color{red} #1}}

\title{Can we `stop' global warming?}
\author{Abba, Baba, Caba, Daba}

\begin{document}

\maketitle

\begin{abstract}
	Something something summarise the article in 300 words or fewer.
\end{abstract}

\vfill
\tableofcontents
\clearpage

\section{Introduction}
Climate change is one of the most significant global environmental challenges,
driven largely by anthropogenic greenhouse gas emissions from fossil fuel combustion,
industrial processes and land-use changes \parencite{bongaarts_climate_2024}.
The resulting emission of carbon dioxide (\cotwo{}) into the atmosphere has caused a sustained rise in global mean temperature,
leading to widespread environmental challenges such as ecosystem degradation,
extreme weather events and sea-level rise.
As part of the initiative of the Global Institute for Climate Change and the Environment (GICCE),
this report seeks to improve general understanding of the Earth’s climate
and develop a simple mathematical model which reflects the goal of international governments to `stop' climate change.

Human activities are the primary driver of increasing concentrations of \cotwo{} in the atmosphere \parencite{friedlingstein_global_2026}.
Fossil fuel combustion, land-use changes and other industrial processes release large quantities of carbon
which exceed the Earth's natural ability to absorb it through carbon sinks \parencite{friedlingstein_global_2026}.
Emissions of \cotwo{} are determined by four key factors, being population, GDP per capita,
energy intensity and carbon intensity \parencite{deutch_decoupling_2017}.
Although technological progress has reduced the carbon intensity of energy by
approximately 0.7\% per year over the past decade,
these improvements have been outweighed by increasing global energy demand as a result
of the growing population \parencite{friedlingstein_global_2026}.
Global anthropogenic \cotwo{} emissions have therefore continued to increase by 0.3\% per year in the decade 2015--2024.
Encouragingly, however, fossil fuel emissions have decreased in 35 economies during the decade 2015--2024
and emissions have slowed from 1.9\% in the decade 2005--2014,
demonstrating the success of partial decoupling of economic growth from emissions through technological changes.

Increased emissions of \cotwo{} into the atmosphere contribute substantially to higher concentrations of greenhouse gases
which trap infrared radiation and thereby raise the global mean temperature.
Global warming is roughly proportional to the cumulative sum of annual carbon emissions \parencite{matthews_proportionality_2009},
the constant of proportionality being the Carbon-Climate Response (CCR).
This relationship forms the basis of the Global Carbon Budget 2025\supercite{friedlingstein_global_2026} (GCB),
which aims to reduce cumulative emissions to limit long-term warming.
The GCB estimates that each additional emission of 180 Gt\cotwo{} will lead to around 0.1\degree{C} of warming,
leaving 170 Gt\cotwo{} remaining from the beginning of 2026 for a 50\% chance of limiting warming to 1.5\degree{C},
which is equivalent to roughly four years of emissions at the current rate \parencite{friedlingstein_global_2026}.

The rising global mean temperature has created significant environmental consequences,
the particular focus of this report being sea-level rises.
Sea-level rise continues even after temperature stabilisation due to the delayed response of oceans and melting ice sheets.
The rate of sea level rise is approximately proportional to global temperature change over long time scales,
so sea level rise can be modelled empirically from temperature \parencite{rahmstorf_semi-empirical_2007}.
If warming is limited to 1.5\degree{C} above the pre-industrial baseline,
substantial long-term sea level rise is still expected,
and is projected to rise approximately 2--3 metres over the next 2,000 years,
although there is low confidence in the exactness of these projections \parencite{friedlingstein_global_2026}.
Reducing future emissions has become increasingly important to achieve temperature stabilisation
and minimise long term environmental impacts.


\section{Literature Review}

The mathematical model produced in this report will quantify interactions between human activity
and the climate system to assess future climate scenarios.
The following literature will be considered according to how effectively each model captures
the relationships required for this investigation and how appropriate their assumptions are for our simple model.

\textcite{matthews_proportionality_2009} related \cotwo{} emissions to global temperature.
Using coupled climate-carbon cycle models, the authors found a linearly proportional relationship
between global warming and cumulative carbon emissions from semi-empirical sources.
The CCR is estimated to be approximately 1--2.1\degree{C} of warming per trillion tonnes of carbon emitted\supercite{matthews_proportionality_2009}.
This relationship is independent of time and atmospheric concentrations.
Further, the model captures the processes of the carbon cycle,
and offsets the effects of carbon uptake into the atmosphere on temperature against carbon sink saturation.
This makes cumulative emissions highly useful for examining scenarios in which temperature is required to stabilise.
However, \citeauthor{matthews_proportionality_2009} calculated that there is substantial uncertainty on the model,
being 0.4--1.5 Tt\cotwo{}.
This is due to the various assumptions on which the relationship relies.
For example, the authors assume that other greenhouse gases and non-greenhouse gases contribute negligibly to global warming.
Further, estimates for historic land-use changes and aerosol forcing from limited available data is assumed to be accurate.
However, uncertainty in this data would have great effects on the CCR value.
The authors also note that their results are derived from idealised climate-carbon model scenarios over decades to centuries,
meaning that extrapolation over longer time scales may not produce accurate results.
The model produced in this report will only rely on \citeauthor{matthews_proportionality_2009}’s linear relationship,
and will estimate a constant of proportionality which is solely based on empirical data.

The Kaya Identity provides a complementary expression for determining annual \cotwo{} emissions.
The identity decomposes annual \cotwo{} into a product of population, GDP per capita, energy intensity and carbon intensity.
\textcite{deutch_decoupling_2017} uses this expression to examine the relationship between economic growth and emissions,
and demonstrates that improvements in energy and carbon intensity can weaken the relationship between economic growth and \cotwo{} emissions.
This is important for the purposes of the current report as it allows a framework for proposing technological changes
to produce a reduction in annual emissions.
However, the Kaya Identity is not a causal model,
so it cannot predict how the terms in the model will respond to technological development,
government policy or changes in the markets.
So, the reliability of the projections in that article are questionable.
This report aims to adopt the relationship established by the Kaya Identity.
However, it will assume that technological advancements have a direct impact on energy and carbon intensity.
The GCB provides empirical critiques against the assumptions of the model.
\textcite{friedlingstein_global_2026} states that global emissions have not yet reached sustained decline
despite energy and carbon intensity improvements,
but also identifies evidence of partial decoupling.
Our model must be closely examined against observed trends.

Population represents a fundamental driver of emissions due to rising demands for energy.
\textcite{lee_demographic_2003} summarised the demographic transition theory,
which describes the growth changes of a population as a region develops.
The paper is helpful to produce a logistics model for long-term demographic projections.
However, the literature identifies several factors which would affect population trajectories
and sheds important authority that a carrying capacity in a simple logistic model may not accurately represent
the social, economic demographic mechanisms which would affect future population.
Our logistic model is therefore used as an approximation for long-term behaviour
and assumes that such behaviour is stable over long timescales.

Stabilising global temperature does not capture all risks associated with climate change.
Sea-level rise continues long after temperatures stabilise due to delayed responses of ocean heat uptake and ice-sheet dynamics \draftmsg{(Source needed, probably Rahmstorf)}.
\textcite{rahmstorf_semi-empirical_2007} proposed a semi-empirical linear relationship between sea-level changes and temperature.
The model captures long-term sea level behaviour as opposed to an assortment of underlying physical processes.
While the model performs well over longer timescales, it is less reliable over shorter periods
as it does not account for fluctuations in physical processes.
Nonetheless, it provides an appropriate framework for assessing whether sea-levels can eventually stabilise
under changing climate scenarios.

The literature provides key relationships required for the model in this report,
which will be applied to three scenarios by the GICCE.
The first is stabilisation of global temperatures to 4\degree{C} above pre-industrial levels by 2100.
The second is stabilisation at 2\degree{C} by 2050,
and the third is limiting warming to no higher than 1.5\degree{C} while plateauing sea-level rise by 2100.
The model is subsequently used to assess the various changes and advancements required to achieve each scenario
and evaluate the implications for the environment and humanity.


\section{Modelling Methodology}

\subsection{Main Models}

To model temperature change, we modified the Transient Climate Response to
Cumulative Carbon Emissions (TCRE) relationship\supercite{matthews_proportionality_2009}.
\citeauthor{matthews_proportionality_2009} established a linear relationship between Global Surface Temperature anomaly
and Cumulative Carbon Emissions done by humans, both since pre industrial levels, which we model as
\begin{equation}
	\Delta T(t) \propto \int_{t_0}^t E(y) dy,	\label{eq:temp_emi}
\end{equation}
where:
\begingroup
\renewcommand{\labelitemi}{-}
\begin{itemize}
    \setlength\itemsep{-0.2ex}
	\item $t$ is time in years;
	\item $t_0$ is the year we start measuring from;
	\item $\Delta T(t)$ is the temperature change at time $t$ since $t_0$; and
	\item $E(t)$ is the amount of \cotwo{} produced by humans in year $t$ measured in in tonnes.
\end{itemize}
\endgroup
% where $\Delta T(t)$ represents temperature at time $t$, and $E$ represents annual emissions produced by humans.
Note that emissions and \cotwo{} emissions are used interchangeably, as \cotwo{} accounts for a majority
of the global heating due to greenhouse gas emissions,
and a majority of greenhouse gases emitted \parencite{noaa_annual_2025}.

We modify this by introducing a reduction term that subtracts from the emissions,
representing human/technological efforts to directly remove carbon from the atmosphere.
This leads to a new equation, which is the primary model for the paper
\begin{equation}
	\Delta T(t) \propto \int_{t_0}^t (E(y) - R(y)) dy,	\label{eq:temp_emi_rem}
\end{equation}
where $R$ is a function representing annual carbon reduction efforts.
To model human carbon emissions $E$, we used the Kaya Identity \parencite{deutch_decoupling_2017}, which states
\begin{equation*}
	E(t) = \text{Population} \times \frac{\text{GDP}}{\text{Population}} \times \frac{\text{Energy used}}{\text{GDP}} \times \frac{\text{\cotwo{} produced}}{\text{Energy used}}.
\end{equation*}
In other words, \cotwo{} emissions is the product of population, GDP per capita, energy intensity, and carbon intensity.
To simplify our model, we capture the product of the latter two as technological inefficiency $D$.
This gives a simplified Kaya Identity
\begin{equation}
	E(t) = P(t) \cdot Y(t) \cdot D(t),	\label{eq:emi_pyd}
\end{equation}
where:
\begingroup
\renewcommand{\labelitemi}{-}
\begin{itemize}
    \setlength\itemsep{-0.2ex}
	\item $P$ is population;
	\item $Y$ denotes GDP per capita; and
	\item $D$ is the technological inefficiency term.
\end{itemize}
\endgroup

Next, we look at these terms in more detail.

\subsection{Sub Models}

\subsubsection{Population}

\subsubsection{Economy}

\subsubsection{Technological Inefficiency}


\section{Results}


\section{Discussion}


\section{Conclusion}


\section{Bibliography}

\begingroup
\nocite{*}
\renewcommand{\section}[2]{}	% Hide the "References" heading
\printbibliography
\endgroup

\section{Appendix}

\end{document}
